% Appendix G -- Messages (PUB-REF-0001 AN-03 order: the message exactly as
% displayed, bold, at the head margin; Explanation:; Action:).  Part 1 is
% the TPS Report Generator's messages in the order the program can display
% them: every message is printed from the table tps_inventory.c generates
% from the locked spec/tps_ritual.json, and test-pubs-tps proves that every
% message of the inventory has its entry and that no entry names a message
% the program lacks (GC-01/GC-02).  The R2 and R5 prompts are the wording
% lane tpsr2 shipped (77835faa, bead initech-5igf).  Part 2 is
% the program's other lines (banner, end of report), quoted from
% os/milton/tpsrep_program.asm m_banner / m_done1 / m_done2.  Part 3 is the
% InitechDOS device error messages, GENERATED from spec/dos_messages.json
% by tools/pubs_msgtable.c (the `device` selector), with the same
% explanations as the InitechDOS Reference Appendix A (os/milton/int21.c;
% 0006 and 0008 are catalogue entries no code path of this version
% displays: DRAFT, bead initech-vlhp); a missing or stray detail fails the
% build (CP-06).
% GATES: test-tps-report test-flair-tps-journey test-pubs-tps test-dosmsg
\Gates{test-tps-report, test-flair-tps-journey, test-pubs-tps, test-dosmsg}
\chapter{Messages}

\section{How to Use This Appendix}
Part 1 lists the messages of the TPS Report Generator, in the order in which
the program can display them.  Each message is shown exactly as it is
displayed, followed by an explanation and the action to take.  Where a
message contains \ph{n}, the program displays a name or a number in its
place.  Part 2 lists the other lines the program displays.  Part 3 lists the
InitechDOS device error messages.

\section{Part 1. Messages of the TPS Report Generator}

\TpsMessageEntry{TPS-001}{The parameter record \file{TPSPARM.DAT} is not in
the current directory.  The program ends with return code 1.}{If the current
directory is not \file{A:\bs APPS}, type \type{CD APPS} and start again at
\proc{02}.  Otherwise see your supervisor.}

\TpsMessageEntry{TPS-002}{The parameter record has no line \ph{n}, where
\ph{n} is \texttt{TITLE}, \texttt{DIVISION}, \texttt{MEMO} or
\texttt{LINES}, the first missing in that order.  The program ends with
return code 1.}{See your supervisor.  The record is corrected under
Operating Procedure OP-7.3.}

\TpsMessageEntry{TPS-010}{The \texttt{LINES} line of the parameter record
\file{TPSPARM.DAT} does not hold a number.  The program ends with return
code 1, before the alignment is asked for and before any file is
changed.}{See your supervisor.  The record is corrected under Operating
Procedure OP-7.3.}

\TpsMessageEntry{TPS-007}{No table was named on the command line, or the
table named is not in the current directory.  The program ends with return
code 7.}{Type the command again with the table name: \type{TPSREP TPS.DBF}.
If the message is displayed again at \texttt{A:\bs APPS>}, see your
supervisor.}

\TpsMessageEntry{TPS-006}{The file named is not a TPS table: \ph{n} is the
first of the six required fields it lacks, or \texttt{SIGNATURE} when it is
not a dBASE III table.  The program ends with return code 6.}{Check the
table name.  If it is \file{TPS.DBF}, see your supervisor.}

\TpsMessageEntry{TPS-011}{The \texttt{PAGES} field of a record of the table
does not hold a number: \ph{n} is the number of the record, counting from
1.  The program ends with return code 6, before the alignment is asked for
and before any file is changed.}{See your supervisor.  Do not correct the
table yourself.}

\TpsMessageEntry{R2.prompt}{The program requires
the form to be aligned.}{Carry out \proc{03}.}

\TpsMessageEntry{R2.refuse}{The alignment attempt was not accepted: \ph{n}
is the number of presses counted.  The count starts again from zero.}{Begin
again at \proc{03}, step 1.}

\TpsMessageEntry{R2.accept}{The form is aligned.}{Continue with \proc{04}.}

\TpsMessageEntry{R3.prompt}{The program asks for the report number, and then
for it again.}{Carry out \proc{04}.}

\TpsMessageEntry{TPS-003}{The two entries of the report number do not
agree, or are empty.  Both are asked for again.}{Begin again at \proc{04},
step 1.}

\TpsMessageEntry{R4.prompt}{The program asks for the number of the memo
about the new cover sheets.}{Carry out \proc{05}.}

\TpsMessageEntry{TPS-004}{The number keyed is not the number in the
\texttt{MEMO=} line of the parameter record.  The program asks
again.}{Check the number on the memo and key it again.  If it is refused a
second time, see your supervisor.}

\TpsMessageEntry{R5.prompt}{The
program asks for the cover sheet attestation.}{Carry out \proc{06}.}

\TpsMessageEntry{TPS-005}{The answer to the cover sheet prompt was N: \ph{n}
is the memo number.  No report is written.  The program ends with return
code 5.}{Attach the new cover sheet, \form{1} (Rev.\ 10/89), and run the
program again from \proc{02}.  If you have no new cover sheet, see your
supervisor.}

\TpsMessageEntry{TPS-008}{\rpt\ could not be created in the current
directory.  No report is written.  The program ends with return code
8.}{See your supervisor.}

\TpsMessageEntry{TPS-009}{\rpt\ could not be written to the disk, usually
because the disk is full.  The report on the screen is not complete and
no report is kept: the unfinished \rpt\ is deleted and the previous report
remains as \file{TPSREPT.BAK}.  The program ends with return code
9.}{See your supervisor.  Do not carry out \proc{07}.}

\TpsMessageEntry{END}{The last line of every run: \ph{n} is the return code
(Session~1).}{If the return code is 0, carry out \proc{07}.  Otherwise carry
out \proc{08}.}

\section{Part 2. Other Lines Displayed by the Program}
\OtherMessage{INITECH CORPORATION}{banner}{The first of the three banner
lines, with \texttt{TPS REPORT GENERATOR\ \ VERSION 2.11} and the copyright
line.  It is displayed at every start.}{Read the line after the banner
(\proc{02}, step 2).}
\OtherMessage{TPS REPORT WRITTEN TO TPSREPT.TXT.}{end of report}{The report
has been written to \rpt.}{Carry out \proc{07}.}
\OtherMessage{PREVIOUS REPORT RETAINED AS TPSREPT.BAK.}{end of report}{The
report that was in \rpt\ before the run is now \file{TPSREPT.BAK}.  It is
displayed only when there was one.}{None.  It is retained until the next
run.}

\section{Part 3. InitechDOS Device Error Messages}
\MessageDetailsOn
These messages are displayed by InitechDOS, not by the program.  The letter
\ph{d} stands for the drive letter.  The number beside each message
identifies it when you call the Help Desk (extension 2504).

\MessageDetail{MSG-DOS-0001}{InitechDOS could not read or write a disk and
asks what to do.  It is displayed after a device error message.}{Correct the
cause and press \type{R}.  If the message is displayed again, press
\type{A} and carry out \proc{08}.}
\MessageDetail{MSG-DOS-0007}{InitechDOS could not read the disk in drive
\ph{d}.}{Press \type{A} and see your supervisor.}
\MessageDetail{MSG-DOS-0006}{The drive \ph{d} was not ready.\draft{see bead
initech-vlhp: catalogue entry not displayed by any code path of this
version}}{Make the drive ready and press \type{R}.}
\MessageDetail{MSG-DOS-0008}{InitechDOS tried to write on a write-protected
disk in drive \ph{d}.\draft{see bead initech-vlhp: catalogue entry not
displayed by any code path of this version}}{Press \type{A} and see your
supervisor.}
\input{dos_messages_device}
\CheckMessageDetails
